\documentclass[12pt,a4paper]{article}

% Norsk språk og tegnsett
\usepackage[T1]{fontenc}
\usepackage[utf8]{inputenc}
\usepackage[norsk]{babel}
\usepackage{todonotes}
\usepackage{tcolorbox}

% Sideoppsett
\usepackage[a4paper, margin=2.5cm]{geometry}
\usepackage{array}
\usepackage{graphicx}
\IfFileExists{phvr8t.tfm}{\usepackage[scaled=0.95]{helvet}}{}
\usepackage{amsmath}

% Referanser
\usepackage[
    style=apa,
    backend=biber
]{biblatex}
\addbibresource{references.bib}

% Lenker
\usepackage[colorlinks=true, allcolors=blue]{hyperref}

% Fyll inn kandidatnummer og ordantall før innlevering.
\newcommand{\kandidatnummer}{[Fyll inn]}
\newcommand{\antallord}{[Fyll inn]}
% Emnenavn: https://www.uio.no/studier/emner/uv/ils/PPU3210/index.html
\newcommand{\emnenavn}{PPU del 1 av 2: Undervisning og læringsforløp}

\begin{document}

% Oppgavetekst på side 0; selve besvarelsen begynner på side 1.
\setcounter{page}{0}
\begingroup
\small
\setlength{\parindent}{0pt}
\setlength{\parskip}{0.45\baselineskip}

{\large\bfseries Semesteroppgave (PPU3210, PPU3510D, PROF3025)\par}

Semesteroppgaven er en individuell skriftlig oppgave som integrerer pedagogikk og fagdidaktikk A. I oppgaven skal du ta utgangspunkt i et undervisningsopplegg du har gjennomført i praksis. Du skal formulere en problemstilling som kan belyse utvalgte sider av dette, og analysere og drøfte undervisningsopplegget ut fra pedagogisk og fagdidaktisk teori.

Problemstillingen skal relateres til ett eller flere av læringsutbyttene beskrevet i emnebeskrivelsen.

\textbf{Formelle krav}

Eksamensbesvarelsen skal ha et omfang på 3000--4000 ord (ikke medregnet forside, litteraturliste og vedlegg). Bruk siste versjon av APA for referanser og litteraturliste.

Det er ikke anledning til å innhente video-, bilde- eller lyddata i forbindelse med semesteroppgaven. Dataene du samler inn (observasjonsnotater etc.), må lagres trygt, som anonymiserte data. Det betyr for eksempel at du ikke skal bruke ekte navn på informanter og på skole, heller ikke i filnavn. All informasjon fra praksis skal anonymiseres i selve oppgaven slik at det ikke er mulig å identifisere skole, lærere, elever og medstudenter. Dette gjelder både for semesteroppgaven og deloppgavene. Undersøkelsen kan ikke offentliggjøres eller brukes i andre sammenhenger enn eksamensbesvarelsen.

Besvarelsen skal ha følgende oppbygning:
\begin{itemize}
    \setlength{\itemsep}{0pt}
    \setlength{\parsep}{0pt}
    \item Forside med kandidatnummer, emnekode og emnenavn, aktuelt semester og antall ord i oppgaven
    \item Innledning med presentasjon av problemstillingen
    \item Kort presentasjon av undervisningsopplegget, inkludert beskrivelse av gruppen/klassen
    \item Analyse og drøfting av undervisningsopplegget ut fra pedagogisk og fagdidaktisk teori, samt erfaringer fra praksis
    \item Begrunnede forslag til hvordan undervisningsopplegget kan videreutvikles
    \item Konklusjon
    \item Litteraturliste
    \item Vedlegg: Undervisningsplan for opplegget. Bruk malen «\href{https://www.uv.uio.no/ils/studier/praksis/undervisningsplan.html}{Undervisningsplan}».
\end{itemize}

Ettfagsstudenter på lektorprogrammet (PROF3025) skal også legge ved liste over selvvalgt tilleggslitteratur. Tilleggslitteraturen må være godkjent av faglærer.

Ettfagsstudenter på PPU deltid (PPU3510D) og PPU heltid (PPU3210) skal bruke ordinært pensum i sitt fag.

Kilde: \href{https://www.uv.uio.no/ils/studier/eksamen/semesteroppgaven-bokmal.html?utm_source=chatgpt.com}{UiO -- Semesteroppgaven (bokmål)}.

\clearpage
\endgroup
\setcounter{page}{1}

% Egen forside på side 1, etter oppgaveteksten på side 0.
\begingroup
\thispagestyle{plain}
\centering
\vspace*{25mm}
{\large Universitetet i Oslo\par}
\vspace{18mm}
{\Huge\bfseries Semesteroppgave\par}
\vspace{10mm}
{\Large \emnenavn\par}
\vspace{25mm}
{\renewcommand{\arraystretch}{1.7}
\begin{tabular}{@{}l@{\hspace{12mm}}l@{}}
    \textbf{Kandidatnummer:} & \kandidatnummer \\
    \textbf{Emnekode:} & PPU3210 \\
    \textbf{Semester:} & Høst 2026 \\
    \textbf{Antall ord:} & \antallord
\end{tabular}\par}
\clearpage
\endgroup

\section{Innledning og presentasjon av problemstillingen}
\label{sec:innledning}

I denne oppgaven så blir vi presentert en undervisningssituasjon innen T-matematikk for en klasse på 11 stk i yrkesfag i VG1 Havbruk blå. Generelt sett er yrkesfag kjent for å ha mindre matematisk modenhet \parencite{lynum2015matte}, men denne klassen er et unntak. 

Da inntaket for å komme inn på denne skolen var godt over 5 i snitt fra 10.klasse, så var dette å anse som meget ambisiøse videregående elever der flere har snakket om å bli veterinær, jobbe innen oppdrettsnæringen eller som fisker. Tradisjonelt er
disse yrkene kjent sett et mannlig yrke, men faktisk i denne klassen var vi tydelig dominert av jenter, der bare 3 av 11 elever var gutt. Alle i klassen er har også bevisst valgt T-matematikk istedenfor P-matematikk. Noen har nevnt at dette er på
grunn av nysgjerrigheten, og andre ser det som nødvendighet for videre utdanning. En enighet er at alle elevene tenker dette faget kan nyttig i fremtiden.  

Vi skal gjennom denne oppgaven prøve å presse denne gjengen litt ekstra, å se om de klarer å håndtere vanskeligere abstrakte oppgaver over tid, uten at det utvikler seg hjelpløshet.

Da spesifikt må man ta hensyn til

























Hvordan kan jeg ivareta alle elevenes læringsbehov og bygge deres utholdenhet og evne til å stå i arbeidet med en vanskelig oppgave over tid. Dette er da en forskjell til T-matematikken  
Kan man  ta hensyn til svake og sterke elever?
             Vi hadde en sterk elev, som faktisk var meget mottagelig for vanskelige oppgaver. Der jeg allerede tidlig tidlig i undervisningen samlet et dokument med eksamensoppgaver innen Vi hadde også en svak elev, som generelt hadde alminnelige regnevanskerm men som hadde ambisiøse planer om å bli veterinær.
Her bant annet vil vi diskutere klasseledelse.




Vi skal også diskutere hvordan vi kan ta være på sterke elever, ved bruk av spesialkunnskap.
Men også hvordan vi bruker horisontal kunnskap til å se hvorfor dette er nyttig og deres kunnskap


Vi skal i dette prosjektet også snakke om hvordan vi går frem i form av rasjonell versus instrumentell forståelse. 

- Koble til Vygotsky med aktuell sone og nærmeste utviklingsone ( hjelp alene versus med læring.)
- Se samling med piatre
- Den der artikkelen med de svake



Ekstra problemstilling:
Kan matematisk modenhet utvikles hos de svake elevene når de introduseres for mer abstrakte og krevende ooppgaver
tidligere i undervisningen, eller
bør de de først få oppleve mestring gjennom enkle oppgaver gradvis og bygge opp et grunnlag med mengdetrening.


\begin{tcolorbox}
    
\label{sec:problemstilling_til_semesteroppgave}


Problemstilling til semesteroppgave:
\begin{itemize}
    \item
Hvordan kan jeg ivareta alle elevenes læringsbehov og bygge deres utholdenhet og evne til å stå i arbeidet med en
vanskelig oppgave over tid.

\item Kan man ta hensyn til svake og sterke elever?
    \begin{itemize}
        \item Koble til klasseledelse
        \item Den sirkelen med læringsbehov ved å bruke horisontal kunnskap, spesialkunnskap, elevens kunnskap etc etc
        \item  Læringskurve: - rasjonell versus instrumentell forståelse.

            - Koble til Vygotsky med aktuell sone og nærmeste utviklingsone ( hjelp alene versus med læring.)
    \end{itemize}
\end{itemize}





- Se samling med piastre
- Den der artikkelen med de svake



Ekstra problemstilling:
Kan matematisk modenhet utvikles hos de svake elevene når de introduseres for mer abstrakte og krevende ooppgaver
tidligere i undervisningen, eller
bør de de først få oppleve mestring gjennom enkle oppgaver gradvis og bygge opp et grunnlag med mengdetrening.
    
\end{tcolorbox}

\section{Presentasjon av undervisningsopplegget}
\label{sec:undervisningsopplegget}


\subsection*{Kort beskrivelse av timen}%
\label{sub:kort_beskrivelse}

% Bruke og forstå kvadratsetningene $(a+ b)^2 = a^2 + 2ab + b^2$. Bli
%   introdusert for egenskapene til andregradsuttrykk i ligninger. Elevene skal også bli kjent med
%   begrepet faktorisering av uttrykk og hvordan dette henger
% sammen med produktregelen, der $ a\cdot b = 0 \implies a = 0 \textit{ eller } b=0 $. Eleven skal
% også kunne faktorisere enkle andregradsuttrykk.

% Elevene har brukt mange undervisningstimer på å utlede den første og
%   den andre kvadratsetningen. De har nå god trening i å gå fra faktorisert form til et utvidet
%   uttrykk. I denne dobbeltimen skal vi prøve å lære elevene hvordan de kan gå fra et utvidet uttrykk
%   til faktorisert form. For å motivere elevene bruker vi produktregelen og ligninger som
%   utgangspunkt, slik at elevene må faktorisere uttrykket for å finne svaret. En av fallgruvene er at
%   produktregelen, som algebraisk sett er et enkelt resultat, kan virke svært abstrakt når vi
%   kombinerer den med faktorisering av andregradsuttrykk.

% For å lage en myk start, så prøvde vi veldig lett å introdusere konseptuelt hva det vil si å dele et tall eller en ukjent på seg selv.
% I dette tilfellet var progresjonen å forenkle
% \[
% \begin{aligned}
%     \text{1a)}\quad &\frac{4x}{x}  ,
%     \qquad \text{1b)}\quad \frac{x(x+1)}{(x+1)} , \\
% \end{aligned}
% \]

% Deretter ga læreren en gul lapp, der de selv skulle ved hjelp av de de har lært regne ut  
% \begin{equation}
%     \text{1c)}\quad &\frac{x^2 + x}{x+1}  \\
% \end{equation}

% Elevene ble da satt til å løse oppgaven på såkalt vertikal tavle, to og to:
% \begin{equation}
%     \text{1d) } \quad  \frac{x^2 + 2x + 1}{x+1}
% \end{equation}

%  I andre time gikk vi frem på samme måte for å introdusere produktregelen, med følgende progresjon:

% \[
% \begin{aligned}
%     \text{2a)}\quad & 4x = 0,
%     \qquad \text{2b)}\quad 4(x-1) = 0, \\
%     \text{2c)}\quad &(x-1)(x-1)= 0
% \end{aligned}
% \]

% Etterpå skulle elevene jobbe alene med tilsvarende oppgaver i boka.
    

\subsection*{Hva gjorde læreren?}%
\label{sub:hva_gjore_laereren}



\subsection*{Hva gjorde elevene?}%
\label{sub:hva_gjorde_elevene}



\subsection*{Hvilken forventet motivasjon og kunnskap hadde elevene}%
\label{sub:elevenes_forutsetninger}

I denne klassen har elevene variert bakgrunn. De kommer fra ulike
  ungdomskoler og sosiale livssituasjoner og har erfaring med ulike lærere og ulik
  undervisningspraksis. Vi kan likevel forvente at de er kjent med ligninger og
regneoperasjoner, men det vil antakelig være behov for oppfriskning og oppklaring av misforståelser.
Det kan også være variasjon i disiplin og utholdenhet, da mange kan ha vansker med å holde
oppmerksomheten over tid når ting blir vanskelige. Det vil si at det kan
være svakere elever i rommet som både har lav utholdenhet og sliter med grunnleggende regneferdigheter.

 



\section{Analyse og drøfting av undervisningsopplegget}
\label{sec:analyse}

\subsection{Erfaringer fra praksis}

%Observert flere ganger i videoen, ser vi elevene være ukonsentrerte. For eksempel ble det tydelig observert at den ene gruppa begynte å kaste kritt på hverandre, noen sover i timen, og på et tidspunkt ble det spilt musikk på toppen av et skap midt i undervisningen. Ut fra videoen kan det også oppleves at mange elever synes matematikk er et vanskelig fag.

%%Som et tiltak for å forbedre klassemiljøet har hun laget en kontrakt som sier at elevene skal følge med i timen og utnytte tiden i skolen for læring, men som åpenbart ikke blir fulgt av elevene. Det skal dog sies at noen elever syntes dette er
%forstyrrende, blant annet da det ble spilt musikk i undervisningen. En annen elev ble faktisk så frustrert over situasjonen og sa ``Kan vi gjøre dette i en annen time, nå har vi matte''. Dette er et viktig poeng, for det innebærer at det er elever i klassen som faktisk ønsker å lære og ikke bare leke.
%
%Gjennom bruk av klassekontrakten har læreren åpenbart brukt mye tid og energi på å tenke på hvordan klassemiljøet kan forbedres. Det er ingen tvil om at når hun først får oppmerksomheten og trigger nysgjerrigheten til elevene, så er de i stand til å lytte og lære. Da hun spurte en elev om hva arealet ble, og brukte litt tid på å tenke på riktig svar, begynte klassen å le litt.
%
%Det blir også observert at mens den ene eleven sover, sier hun at eleven ikke kan være trøtt når man gjør matematikk. Dette er verdt å merke seg, da søvn henger dypt sammen med fysiologiske behov for optimal læring. 
%
%Det er generelt relativt lite dialog mellom elev og lærer i undervisningen. For å komme til ordet i timen beordrer hun oppmerksomhet med å hysje på elever. Hun øker også strengheten ved at hvis ikke elever rekker opp hånda under felles diskusjon, så får eleven anmerkning. Så er spørsmålet om denne autoritære måten å håndtere elevene på vil hjelpe dem å arbeide i faget.

\section{Begrunnede forslag til forbedring av undervisningen}
\label{sec:forbedringer}

\subsection{Relasjoner og klassemiljø}

%Et annet viktig spørsmål er hvordan relasjonen mellom lærer og elev påvirker motivasjonen. Læreren hysjer flere ganger på elevene, og elevene spøker samtidig med læreren og spiller musikk i undervisningen. Dette kan tyde på at læreren har utfordringer med å etablere tydelige rammer for undervisningen. I den andre videoen sover en elev, og læreren sier på en rolig måte at han må henge opp jakken. Det kan være flere grunner til at eleven sover i timen. Det kan for eksempel være mangel på søvn, noe som ikke nødvendigvis betyr at eleven ikke ønsker å lære. I stedet for å tolke atferden som manglende motivasjon, kan læreren undersøke hva som ligger bak. 

% Her må læreren gå virkelig til verks for å bli kjent med elevene både i og utenfor klasserommet. Hvis en elev ofte sover i timen, kan læreren ta en prat med eleven og spørre om hva han gjør på kvelden, hvordan han sover, og om det er noe som gjør det vanskelig å møte uthvilt på skolen. 
% Hvis man følger Maslows behovspyramide, vil man først søke å dekke fysiologiske behov, trygghet og tilhørighet før man frivillig oppsøker intellektuelle behov. Man kan derfor med fordel gå tidlig til verks dersom det er noe åpenbart man kan løse
% nederst i pyramiden \parencite{skaalvik2021skolen}. Læreren bør spørre om trivsel, venner, hvordan de opplever skolen og hva de gruer seg til.  Også dette kan sees i syn av Maslows behovspyramide, da bygging av relasjon vil gi anerkjennelse og trygghet i klasserommet som senere er et godt verktøy for videre utvikling.
% 
% Elevene må kunne oppleve at det de sier blir tatt seriøst. Dersom elevene opplever at ideene deres blir møtt med respekt, kan det skape et tryggere læringsmiljø. Dette er særlig viktig i matematikk, hvor elever ofte må prøve seg frem og gjøre feil.
% Forskning har vist at for et godt læremiljø, så er det viktig at elevene viser troverdighet til læreren i form av at deres ideer blir mottatt seriøst og at eventuell diskusjon er basert på faglig begrunnelse og ikke som individ \parencite{ucar}.
% Dette er også viktig for danning av eleven, da dette skaper en demokratisk platform i klasserommet der hver mening skal bli hørt og man sammen skal utvikle egen og hverandres læring. 
% Spesielt er dette essensielt ved elevenes valg om å forstyrre i timen, da man potensielt kan påvirke andre enn seg selv \parencite{werler2024}. I videoen kan man blant annet observere at enkelte elever er ufokuserte og lager støy i klasserommet. Som reaksjon sier en annen elev «nå har vi matte», og da er det viktig at man som demokratisk individ i klassen ikke bestemmer fokuset for resten av klassen.
%   



\subsection{Individuell opplæring}
% Hva er det som får disse elevene til å yte høy innsats i timen? Målorientering handler om at elever har ulike grunner til å arbeide med skolefag, og dermed ulike grunner til å yte høy eller lav innsats. I denne sammenhengen kan vi skille mellom læringsorientering og egoorientering. En læringsorientert elev er opptatt av å utvikle egen kompetanse og lære av feil. Eleven søker optimale utfordringer og ser på feil som noe man kan lære av. Forskning har vist at læringsorientert motivasjon kan bidra til større utholdenhet når eleven møter utfordrende problemer \parencite{skaalvik2026}.
% 
% Ved egoorientering står eleven selv mer i sentrum. Eleven arbeider for å bli positivt vurdert av andre eller for å være blant de beste. Eleven er ute etter å demonstrere kompetanse og legger større vekt på sosial sammenligning. Feil kan derfor
% oppfattes som truende, og eleven kan lettere gi opp når oppgaven blir vanskelig \parencite{skaalvik2026}. Sett fra lys av Covingtons teori om selvverd, så viser det seg at egne evner, gode prestasjoner og personlig suksess har avgjørende betydning
% for selvverd og videre motivasjon. Derfor må undervisningen tilpasses elevenes forutsetninger og ståsted, som potensielt kan være en strategi for å bygge opp motivasjon og selvverd. Her er det også viktig å dempe
% elevenes tendenser til sosial sammenligning konkurranse, da man individuelt stiller på forskjellige plan og utfordringer og ofte er ikke alle vinnere \parencite{skaalvik2021skolen}.   
% 
% Å legge til rette for et godt læringsmiljø, der ingen spørsmål er dumme, er derfor viktig for å fremme nysgjerrighet i klassen. Dersom elevene opplever at det er trygt å komme med forslag og svare feil, kan det bli lettere å delta aktivt i matematikkundervisningen. En læringsorientert kultur kan dermed bidra til at elevene ser feil som en naturlig del av læringen, fremfor noe som bør unngås. Dog bør man være oppmerksom på at spørsmål fra enkeltelever kan skape reaksjoner i klassen. Som
% klasseleder er det derfor viktig å håndtere slike reaksjoner på en måte som ikke bidrar til at det skjer igjen. Dersom slike reaksjoner oppstår, bør de tas opp tidlig for å hindre at det utvikler seg til mobbing, noe som kan ha en negativ innvirkning
% på læringsmiljøet \parencite{werler2024}.



\section{Konklusjon}
\label{sec:konklusjon}

%Det vi ser er at det finnes tiltak som er verdt å prøve for klassen. Man kan først gå til verks å bygge relasjon med elev og lærer, og blant annet identifisere hvert elevs behov iforhold til Maslows behovshiarki muligens se om eleven er i en tilstand
% der læring er mulig. Deretter kan man bruke klasserommet som en helhet å lære elevene å leve i et demokrati i klassen, der hver svar og forstyrrelse vil påvirke alle sammen. Et andre handlingsalternativ er å gi elevene
% individuell opplæring, spesielt da vi kan se i lys fra Covingtons teori om selvverd at å identifisere elevens evner og potensiale for suksess er essensielt for videre motivasjon og selvverd.  
% 

\section{Litteraturliste}
\label{sec:litteraturliste}
% Litteraturlisten fylles automatisk når kildene siteres i teksten.
\printbibliography[heading=none]

% Vedlegg: kopi av den utfylte undervisningsplanen fra arbeidskrav 2.
% Redigerbar undervisningsplan fra arbeidskrav 2.
% Alle tre delene stables på én stående A4-side og skaleres samlet.
\clearpage
\appendix
\newgeometry{left=10mm,right=10mm,top=10mm,bottom=10mm}
\begingroup
\fontsize{10}{12}\selectfont
\sffamily
\setlength{\parindent}{0pt}
\setlength{\parskip}{0pt}
\setlength{\tabcolsep}{1.6mm}
\setlength{\arrayrulewidth}{0.5pt}
\renewcommand{\arraystretch}{1}
\raggedright
\refstepcounter{section}
\label{app:undervisningsplan}
\addcontentsline{toc}{section}{Vedlegg \thesection: Undervisningsplan}
\newlength{\planoriginalbredde}
\setlength{\planoriginalbredde}{274mm}

\ExplSyntaxOn
\prop_new:N \g_plan_prop
\NewDocumentCommand{\fyllinn}{m +m}
  {\prop_gput:Nnn \g_plan_prop {#1} {#2}}
\ExplSyntaxOff
% Grunnopplysninger
\fyllinn{Dato}{19.10.26}
\fyllinn{Trinn}{VG1}
\fyllinn{Fag}{1T}
\fyllinn{Varighet}{2 timer}
% Kompetansemål
% Ett mål per \item.
\fyllinn{Kompetansemål}{%
  \item forklare forskjellen mellom en identitet, en ligning,
        et algebraisk uttrykk og en funksjon
  \item utforske strategier for å løse ligninger, ligningssystemer og
        ulikheter og argumentere for tenkemåtene sine
  \item lese og forstå matematiske bevis og utforske og utvikle bevis
        i relevante matematiske emner
}
% Mål og innhold
% Hva skal elevene kunne, forstå eller mestre etter undervisningen?
\fyllinn{Læringsmål}{

    Eleven skal bli kjent med forkorting av uttrykk av høyere grad.
    Eleven skal forstå hvordan faktorisering og produktregelen kan brukes for å løse ligninger.

}
% Hvilke begreper, temaer og ferdigheter skal undervisningen handle om?
\fyllinn{Faglig innhold}{ Bruke og forstå kvadratsetningene $(a+ b)^2 = a^2 + 2ab + b^2$. Bli
  introdusert for egenskapene til andregradsuttrykk i ligninger. Elevene skal også bli kjent med
  begrepet faktorisering av uttrykk og hvordan dette henger
sammen med produktregelen, der $ a\cdot b = 0 \implies a = 0 \textit{ eller } b=0 $. Eleven skal
også kunne faktorisere enkle andregradsuttrykk. }
% Hvorfor har du valgt disse læringsmålene og dette faglige innholdet?
\fyllinn{Begrunnelse for mål}{ Elevene har brukt mange undervisningstimer på å utlede den første og
  den andre kvadratsetningen. De har nå god trening i å gå fra faktorisert form til et utvidet
  uttrykk. I denne dobbeltimen skal vi prøve å lære elevene hvordan de kan gå fra et utvidet uttrykk
  til faktorisert form. For å motivere elevene bruker vi produktregelen og ligninger som
  utgangspunkt, slik at elevene må faktorisere uttrykket for å finne svaret. En av fallgruvene er at
  produktregelen, som algebraisk sett er et enkelt resultat, kan virke svært abstrakt når vi
  kombinerer den med faktorisering av andregradsuttrykk.
}
% Gruppe- og situasjonsbeskrivelse
\fyllinn{Rammefaktorer}{
    Undervisningen gjennomføres med 12 elever på VG1 i matematikk 1T, i en dobbeltime
    på to ganger 45 minutter med pause mellom timene. Elevene har variert matematisk
    bakgrunn og møter nye og abstrakte begreper. Arbeidet veksler mellom felles
    tavleundervisning, individuelle oppgaver på gule lapper, pararbeid på vertikale
    tavler og oppgaver i læreboka.
}
\fyllinn{Begrunnelse for rammefaktorer}{
    Tidsrammen krever tydelige overganger mellom felles gjennomgang og elevaktivitet.
    Gruppestørrelsen gir læreren mulighet til å følge opp enkeltpersoner og par.
    Variasjonen i elevenes forkunnskaper gjør det nødvendig å følge med på forståelse
    og utholdenhet underveis. Tavler, gule lapper og lærebok må være tilgjengelige
    for de planlagte aktivitetene.
}
% Elevgruppen og vurdering
% Hvilke forkunnskaper har elevene? Hvilke behov for tilpasning ser du?
\fyllinn{Læreforutsetninger}{ I denne klassen har elevene variert bakgrunn. De kommer fra ulike
  ungdomskoler og sosiale livssituasjoner og har erfaring med ulike lærere og ulik
  undervisningspraksis. Vi kan likevel forvente at de er kjent med ligninger og
regneoperasjoner, men det vil antakelig være behov for oppfriskning og oppklaring av misforståelser.
Det kan også være variasjon i disiplin og utholdenhet, da mange kan ha vansker med å holde
oppmerksomheten over tid når ting blir vanskelige. Det vil si at det kan
være svakere elever i rommet som både har lav utholdenhet og sliter med grunnleggende regneferdigheter.   }
% Hvorfor passer tilpasningene til denne elevgruppen?
\fyllinn{Begrunnelse for tilpasning}{

     % Det er verdt å nevne at dette er et fag som viser store krav til tilegnelse av abstrakt informasjon, og elevene har tatt et aktivt valg å velge dette grunnet utdanningsmuligheter, nyskjerrighet og ambisiøsitet.

    For å bygge et godt grunnlag ønsker vi å utfordre elevene til å utvikle generell utholdenhet i
    arbeidet med å tilegne seg ny, abstrakt informasjon. Samtidig ønsker vi å unngå å presse de
    svake elevene så mye at de opplever hjelpesløshet for tidlig i undervisningen.
    Derfor er undervisningen bygget opp ved at vi gradvis legger til byggeklosser med logiske
    premisser. De sterke
    elevene utvikler også forståelsen sin ved å utforske nytten av produktregelen og dens begrensninger.
}
% Hvordan får du informasjon om elevenes læring? Hva skal du se etter?
% Hva skal eventuelt vurderes, og hvordan gir du tilbakemeldinger?
\fyllinn{Vurdering}{ Det ble også stilt muntlige kontrollspørsmål kontinuerlig gjennom
  tavleundervisningen, både til svake og sterke elever. Vi sjekket også utregningene på de gule
  lappene for å fange opp vanlige feil. Under arbeidet på vertikal tavle var
vurderingen basert på elevenes muntlige forklaringer til hverandre.}
% Hvorfor velger du disse vurderingsformene?
\fyllinn{Begrunnelse for vurdering}{ Dette er en effektiv og sosial vurderingsmåte. Da elevene i
  denne undervisningen lærte om nye verktøy i matematikk, det vil si å kombinere faktorisering og
  produktregelen, forventet vi ikke at de skulle forstå dette i løpet av denne timen.
Et av målene var å teste elevenes utholdenhet i arbeidet med å tilegne seg abstrakt informasjon.       } % Gjennomføring
% Når skjer de ulike aktivitetene, og hvor lenge varer de?
\fyllinn{Tidsplan}{%
  \textbf{1. time}\par\smallskip
  10 min\par
  1a)--1b)\par\smallskip
  10 min\par
  1c)\par\smallskip
  25 min\par
  1d)\par\medskip
  \textbf{Pause}\par\medskip
  \textbf{2. time}\par\smallskip
  10 min\par
  2a)--2b)\par\smallskip
  10 min\par
  2c)\par\smallskip
  25 min\par
  Oppgaver\par
  i boka
}
% Hva gjør læreren? Hvordan organiseres undervisningen?
\fyllinn{Lærerens aktiviteter}{

I første time brukte vi enkle eksempler for å illustrere hva vi skulle oppnå.
I dette tilfellet var progresjonen å forenkle

\[
\begin{aligned}
    \text{1a)}\quad &\frac{4x}{x}  ,
    \qquad \text{1b)}\quad \frac{x(x+1)}{(x+1)} , \\
    \text{1c)}\quad &\frac{x^2 + x}{x+1}  \\
\end{aligned}
\]

Elevene ble da satt til å løse oppgaven på såkalt vertikal tavle, to og to:
\begin{equation*}
    \text{1d) }\frac{x^2 + 2x + 1}{x+1}
\end{equation*}

 I andre time gikk vi frem på samme måte for å introdusere produktregelen, med følgende progresjon:

\[
\begin{aligned}
    \text{2a)}\quad & 4x = 0,
    \qquad \text{2b)}\quad 4(x-1) = 0, \\
    \text{2c)}\quad &(x-1)(x-1)= 0
\end{aligned}
\]

Etterpå skulle elevene jobbe alene med tilsvarende oppgaver i boka.
}

% Hva gjør elevene? Arbeider de alene, i grupper eller i fellesskap?
\fyllinn{Elevenes aktiviteter}{I oppgavene 1a) og 1b) valgte vi tradisjonell tavleundervisning, der
  elevene fulgte med på tavla. I oppgave 1c) fikk alle elevene en gul lapp der de selvstendig skulle
  finne et svar, før vi gikk gjennom fasiten på tavla.

}
% Hvilke læremidler og ressurser brukes? F.eks. lærebok og digitale verktøy.
\fyllinn{Læringsressurser}{ Tavleundervisning, vertikal tavle, gul lapp}
% Hvorfor velger du disse arbeidsformene, aktivitetene og ressursene?
\fyllinn{Begrunnelse for aktiviteter}{

Allerede i oppgave $1c)$ mistenkte vi at elevene måtte tilegne seg ny kunnskap, da de måtte kunne
faktorisere $x^2 + x$. Derfor
    valgte vi å gi elevene en gul lapp hver, der de skulle prøve å løse oppgaven i stillhet. Denne
    gule lappen er i prinsippet et verktøy for at de her og nå skal forsøke å løse oppgaven med egne
    ord og tanker.

    Hensikten med vertikal tavle i dette tilfellet er at elevene skal samarbeide om å løse nye
    oppgaver ved hjelp av kreativitet og problemløsning.

    De hadde allerede lært mye nytt i første time, og i mange tilfeller ville en lærer ha stoppet
    opp og repetert det elevene hadde lært. Da vi ville teste utholdenheten deres, valgte vi å
    introdusere produktregelen og faktorisering, som er nye, abstrakte verktøy
    for mange. Dette er viktig for å sette dem i en situasjon der de er nødt til å tenke, selv etter
    å ha lært mye i første time. Senere ble de satt til å løse oppgaver i boka individuelt.

}
% Veiledning
% Hva ønsker du veiledning på? Hvilke spørsmål vil du ta opp?
\fyllinn{Spørsmål til veiledning}{
    I andre time gikk vi bevisst utenfor elevenes komfortsone og satte dem i en situasjon der mange
    i klassen ikke var i stand til å løse problemet. Kan vi potensielt sette for mye press på
    svakere elever?
}
% Felt for veilederens tilbakemeldinger.
\fyllinn{Veileders kommentarer}{

    Generelt sett kan man stille seg følgende spørsmål:
    \begin{itemize}
        \item
Hvordan ivaretar du elevenes behov for fremovermeldinger (vurdering for læring) i undervisningsøkten din?
        \item
Hvordan kan elevene selv ta ansvar for egen læring ved hjelp av læringsmålet for timen?
        \item
Hvordan skal du klare å beholde elevenes fokus gjennom tavleundervisningen?
    \end{itemize}

    Hvis disse spørsmålene er belyst og svarene er godt begrunnet, har du et godt utgangspunkt for å
    analysere den enkelte elevens læringsmiljø.

}

\newcommand{\felt}[1]{{\color{gray}\normalfont\itshape[#1]}}
\ExplSyntaxOn
\NewDocumentCommand{\hent}{m}
  {\prop_item:Nn \g_plan_prop {#1}}
\NewDocumentCommand{\visfelt}{m m}
  {
    \prop_get:NnN \g_plan_prop {#1} \l_tmpa_tl
    \tl_if_blank:VTF \l_tmpa_tl {\felt{#2}} {\tl_use:N \l_tmpa_tl}
  }
\ExplSyntaxOff
\newcolumntype{P}[1]{>{\raggedright\arraybackslash}p{#1}}
\newsavebox{\celleboks}
\newcommand{\mincelle}[2]{%
  \sbox{\celleboks}{\begin{minipage}[t]{\linewidth}\raggedright#2\end{minipage}}%
  \ifdim\dimexpr\ht\celleboks+\dp\celleboks\relax<#1\relax
    \begin{minipage}[t][#1][t]{\linewidth}\raggedright#2\end{minipage}%
  \else\usebox{\celleboks}\fi}
\newcommand{\etikett}[1]{{\bfseries #1}\par}
\newcommand{\sporsmal}[1]{{\itshape #1}\par}
\newlength{\trekolonner}
\setlength{\trekolonner}{\dimexpr\planoriginalbredde-6\tabcolsep-4\arrayrulewidth\relax}
\newlength{\sekskolonner}
\setlength{\sekskolonner}{\dimexpr\planoriginalbredde-12\tabcolsep-7\arrayrulewidth\relax}
\newlength{\firekolonner}
\setlength{\firekolonner}{\dimexpr\planoriginalbredde-8\tabcolsep-5\arrayrulewidth\relax}



% Bygg hele planen som LaTeX-innhold før skalering.
\newsavebox{\planhelhet}
\newsavebox{\planskalert}
\sbox{\planhelhet}{%
  \begin{minipage}{\planoriginalbredde}
% ==================== UNDERVISNINGSPLAN: SIDE 3 ====================
\vspace*{4mm}
{\centering\fontsize{16}{19}\selectfont\bfseries Vedlegg \thesection: Undervisningsplan\par}
\vspace{7mm}

\begin{tabular}{|P{.16\firekolonner}|P{.22\firekolonner}|P{.13\firekolonner}|P{.49\firekolonner}|}
  \hline
  \mincelle{13mm}{\textbf{Dato:} \hent{Dato}} &
  \textbf{Trinn:} \hent{Trinn}                &
  \textbf{Fag:} \hent{Fag}                    &
  \textbf{Varighet:} \hent{Varighet}            \\
  \hline
\end{tabular}

\vspace{6mm}
\textbf{Kompetansemål:}
\vspace{4mm}
\begin{list}{-}{\setlength{\leftmargin}{12mm}\setlength{\labelsep}{4mm}%
    \setlength{\topsep}{0pt}\setlength{\itemsep}{1pt}\setlength{\parsep}{0pt}}
  \hent{Kompetansemål}
\end{list}
\vspace{19mm}

{\fontsize{9}{10.8}\selectfont
  \begin{tabular}{|P{.218\trekolonner}|P{.443\trekolonner}|P{.339\trekolonner}|}
    \hline
    \mincelle{9mm}{\textbf{Beskrivelse}} &
    \textbf{Dine notater}                &
    \textbf{Begrunnelse:} Hvorfor gjør du dette valget? \\
    \hline
    \mincelle{35mm}{\etikett{Mål}
      \sporsmal{Læringsmål: Operasjonalisering av kompetansemål.}
      \vspace{2mm}
      \sporsmal{Kunnskaper, ferdigheter og forståelse som elevene skal utvikle, både faglige og sosiale mål.}
      \vspace{2mm}
      \sporsmal{Faglig innhold.}}
                                         &
    \etikett{Læringsmål}
    \vspace{2mm}
    \visfelt{Læringsmål}{Hva skal elevene kunne, forstå eller mestre etter undervisningen?}
    \par\vspace{4mm}
    \visfelt{Faglig innhold}{Sentrale begreper og faglig innhold}
                                         &
    \visfelt{Begrunnelse for mål}{Begrunn valget av læringsmål og faglig innhold.}
    \\
    \hline
    \mincelle{30mm}{\etikett{Rammefaktorer}
      \sporsmal{Gruppe- og situasjonsbeskrivelse. Hvilke rammefaktorer vil påvirke undervisningen? (f.eks. rom, gruppestørrelse, læringsmiljø, elevgruppens mangfold og (anonymiserte) enkeltelevers utfordringer)}}
                                         &
    \visfelt{Rammefaktorer}{Gruppe- og situasjonsbeskrivelse og praktiske rammer}
                                         &
    \visfelt{Begrunnelse for rammefaktorer}{Begrunn hvordan rammefaktorene påvirker undervisningen.}
    \\
    \hline
    \mincelle{30mm}{\etikett{Læreforutsetninger/Tilpasning}
      \sporsmal{Hvilke forutsetninger har elevene for å lære dette? Hvilke behov for tilpasning ser du?}}
                                         &
    \visfelt{Læreforutsetninger}{Forkunnskaper, elevgruppe og behov for tilpasning}
                                         &
    \visfelt{Begrunnelse for tilpasning}{Begrunn hvordan undervisningen tilpasses elevene.}
    \\
    \hline
  \end{tabular}
}

\par\vspace{6mm}

% ==================== UNDERVISNINGSPLAN: SIDE 4 ====================
{\fontsize{9}{10.8}\selectfont
  \begin{tabular}{|P{.218\trekolonner}|P{.443\trekolonner}|P{.339\trekolonner}|}
    \hline
    \mincelle{29mm}{\etikett{Vurdering og tilbakemelding}
      \sporsmal{Hvordan får du informasjon om elevenes læring? Hva skal du se etter? Hva skal eventuelt vurderes og hvordan?}}
     &
    \visfelt{Vurdering}{Hvordan følger du elevenes læring, og hvordan gir du tilbakemelding?}
     &
    \visfelt{Begrunnelse for vurdering}{Begrunn vurderingsformene og hva de gir informasjon om.}
    \\
    \hline
  \end{tabular}
}

\vspace{4mm}

% Aktivitetstabellen henter innhold fra del 5 øverst i filen.
{\fontsize{9}{10.8}\selectfont
  \begin{tabular}{|P{.185\sekskolonner}|P{.065\sekskolonner}|P{.265\sekskolonner}|P{.12\sekskolonner}|P{.11\sekskolonner}|P{.255\sekskolonner}|}
    \hline
    \mincelle{17mm}{\etikett{Arbeidsformer og aktiviteter. Læremidler og læringsressurser.}}
     & \textbf{Tid}
     & \textbf{Hva gjør læreren?}
     & \textbf{Hva gjør elevene?}
     & \textbf{Læringsressurs:}
     & \textbf{Begrunnelse:}                                                                   \\
    \hline
    \mincelle{45mm}{\sporsmal{Hvem gjør hva? Når? Hvordan vil du organisere elevenes læringsaktiviteter? Hvilke arbeidsformer vil du benytte? (F.eks. presentasjon, individuelt arbeid eller gruppearbeid, utforskende oppgaver, klassesamtale)}
      \vspace{3mm}
      \sporsmal{Hvilke læringsressurser vil du bruke? F.eks. lærebok, digitale ressurser)}}
     & \visfelt{Tidsplan}{Tid}
     & \visfelt{Lærerens aktiviteter}{Lærerens aktiviteter og organisering}
     & \visfelt{Elevenes aktiviteter}{Elevenes aktiviteter}
     & \visfelt{Læringsressurser}{Materiell og ressurser}
     & \visfelt{Begrunnelse for aktiviteter}{Begrunn arbeidsformer, aktiviteter og ressurser.} \\
    \hline
  \end{tabular}
}

\vspace{8mm}

% Veiledningsfeltet har én ytre ramme som i referansedokumentet.
\begingroup
\setlength{\fboxsep}{1.5mm}
\fbox{\begin{minipage}[t]{\dimexpr\textwidth-2\fboxsep-2\fboxrule\relax}
    \mincelle{55mm}{%
      \textbf{TIL BRUK I VEILEDNINGSSITUASJONER}
      {\fontsize{9}{11}\selectfont (rubrikkene kan utvides etter behov)}
      \par\vspace{3mm}
      \begin{minipage}[t]{.49\linewidth}\raggedright
        \textbf{Hva ønsker jeg veiledning på?}
        \par\vspace{5mm}
        \visfelt{Spørsmål til veiledning}{Spørsmål og temaer til veiledning}
      \end{minipage}\hfill
      \begin{minipage}[t]{.49\linewidth}\raggedright
        \textbf{Veileders kommentarer:}
        \par\vspace{5mm}
        \visfelt{Veileders kommentarer}{Kommentarer}
      \end{minipage}
    }
  \end{minipage}}
\endgroup


  \end{minipage}%
}
% Tilpass først bredden, og reduser videre hvis høyden krever det.
\sbox{\planskalert}{\resizebox{\linewidth}{!}{\usebox{\planhelhet}}}
\noindent\makebox[\linewidth][c]{%
  \ifdim\dimexpr\ht\planskalert+\dp\planskalert\relax>0.97\textheight
    \resizebox*{!}{0.97\textheight}{\usebox{\planskalert}}%
  \else
    \usebox{\planskalert}%
  \fi
}\par
\clearpage
\endgroup
\restoregeometry


\end{document}
